% ============================================================
% 55_readability.tex — Intelligente Zeilenumbruch-Steuerung
% ============================================================
% Ziel: In schmalen Spalten (hier 4-spaltiges Landscape-A4) bricht
% der Satz lieber an natürlichen Wortgrenzen, statt die Zeile mit
% Leerraum aufzublähen oder Wörter hart zu trennen, nur um den
% rechten Rand zu erreichen.
%
% Kernbausteine:
%   \ZSFReadableParams : feingetunte TeX-Penalties (Flattersatz-kompatibel)
%   \ZSFReadableOn     : aktiviert RaggedRight + Parameter im aktuellen Scope
%   \ZSFReadableAuto   : respektiert den globalen Schalter
%   ZSFReadable        : Environment für lokale Anwendung
%   \ZSFbreak / \ZSFnobreak / \ZSFallowbreak : semantische Feinsteuerung
%   \ifZSFReadableBody : globaler Schalter (Standard: an)
%
% Dokumentweite Regel: Direkt nach \begin{multicols*} aktiviert
% \ZSFReadableAuto den Flattersatz für den gesamten Spalteninhalt.
% Lokale Anwendungen in Boxen bleiben erlaubt und sind unschädlich.
% ============================================================

% ------------------------------------------------------------
% Globaler Schalter: erlaubt das System projektweit aus-/einzuschalten
% ------------------------------------------------------------
\newif\ifZSFReadableBody
\ZSFReadableBodytrue
\newcommand{\ZSFReadableBodyOn}{\ZSFReadableBodytrue}
\newcommand{\ZSFReadableBodyOff}{\ZSFReadableBodyfalse}

% ------------------------------------------------------------
% TeX-Parameter für "lieber früh brechen als quetschen / hyphen"
% Werte bewusst moderat, damit Silbentrennung möglich bleibt, aber
% nicht die erste Wahl ist (wichtig in 4-Spalten-Layout).
% ------------------------------------------------------------
\newcommand{\ZSFReadableParams}{%
  \pretolerance=1200          % erst ohne Hyphenation versuchen (Default 100)
  \tolerance=2400             % dann Hyphenation erlauben (Default 200)
  \hyphenpenalty=200          % Silbentrennung leicht teurer -> Wortgrenze bevorzugt
  \exhyphenpenalty=200        %
  \doublehyphendemerits=10000 % zwei Bindestriche in Folge vermeiden
  \finalhyphendemerits=10000  % letzter Absatz-Zeile nicht hyphenisieren
  \emergencystretch=\ZSFreadabilityEmergencyStretch % Notfall-Dehnung gegen overfull hbox
  \hbadness=2000              % leise bei kleinen underfulls
  \linepenalty=10             % extra Zeilen sind günstig -> weniger Quetschen
}

% ------------------------------------------------------------
% \ZSFReadableOn: Flattersatz (RaggedRight aus ragged2e) + Parameter
%   - RaggedRight behält Silbentrennung als Fallback bei.
%   - rightskip hat etwas Flexibilität, damit kein "Waterfall"-Look entsteht.
% ------------------------------------------------------------
\newcommand{\ZSFReadableOn}{%
  % Schmale Flatterzone: natürliche Wortabstände, aber gut gefüllte Spalten.
  \setlength{\RaggedRightRightskip}{\ZSFreadabilityRaggedRightskip}%
  \RaggedRight
  \ZSFReadableParams
}

% Environment-Form für lokalen Einsatz
\newenvironment{ZSFReadable}{\par\ZSFReadableOn\ignorespaces}{\par}

% ------------------------------------------------------------
% Komfort-Wrapper: bedingtes Aktivieren (gesteuert über \ifZSFReadableBody)
% ------------------------------------------------------------
\newcommand{\ZSFReadableAuto}{\ifZSFReadableBody\ZSFReadableOn\fi}

%
%
%
%
%
%
%
%
%
%
%
%
%
%
%
%
%
%
%
%
\newcommand{\ZSFbreak}{\unskip\penalty-50\hskip\fontdimen2\font
  plus\fontdimen3\font minus\fontdimen4\font\relax}

% \ZSFnobreak — verhindert Umbruch an dieser Stelle (gebundenes Leerzeichen).
%   Nutzbar für Einheiten: "5\ZSFnobreak kg", "Punktladung\ZSFnobreak q_1"
\newcommand{\ZSFnobreak}{\nobreak\ }

% \ZSFallowbreak — erlaubt Umbruch an einer sonst FESTEN Stelle, also dort, wo
%   kein Leerzeichen steht ("Wortgrenze\ZSFallowbreak Fortsetzung" bleibt ein
%   Wort). Deshalb bringt es bewusst KEINEN Zwischenraum mit — das ist der
%   Unterschied zu \ZSFbreak und nicht eine vergessene Zeile.
%
%   \penalty0 statt \penalty100: »erlaubt« heisst neutral. Eine positive Strafe
%   erlaubt den Umbruch zwar auch, macht ihn aber teurer als jede andere
%   Stelle — der Name sagte damit etwas anderes als der Wert tat.
\newcommand{\ZSFallowbreak}{\penalty0\relax}

%
%
%
%
%
%
%
%
%
%
%
%
%
%
%
%
%
%
%
%
%
%
%
%
%
\newcommand{\ZSFcolonBindPenalty}{9000}
\makeatletter
\newcommand{\ZSF@bind}[1]{%
  \ifmmode~\else\unskip\penalty#1\hskip\fontdimen2\font\relax\fi
}
% Die beiden oeffentlichen Marken MUESSEN in derselben \makeatletter-Klammer
% stehen: \ZSF@bind ist ein Kontrollwort mit @, und nach \makeatother zerfiele
% es beim Tokenisieren in \ZSF und @bind.
\newcommand{\ZSFcolonbind}{\ZSF@bind{\ZSFcolonBindPenalty}}
\newcommand{\ZSFunitbind}{\ZSF@bind{10000}}
\makeatother

% ------------------------------------------------------------
% Automatischer LuaLaTeX-Filter für intelligente Zeilenumbrüche
% ------------------------------------------------------------
% Lädt den Lua-Filter zur Vermeidung von Trennungen bei Doppelpunkten,
% Einheiten und kurzen Inline-Formeln. Die Registrierung erfolgt
% über \AtBeginDocument, damit Paketladevorgänge ungestört bleiben.
\directlua{dofile(kpse.find_file("styles/zsf_filter.lua"))}
\AtBeginDocument{%
  \directlua{zsf_register_filter()}%
}

% ============================================================
